\subsection{Central lifting and the remaining compatibility problem}

The local theorem checked in D.19 has a further global consequence.  It
removes the obstruction to every finite-rank cuspidally central lifting
problem.  This assertion concerns a fixed base curve; passage through
its finite étale neighbourhoods requires additional compatibility.

\begin{claim}{D.23}{Proved}\label{clm:D23}
Let $V/k$ be a smooth geometrically connected hyperbolic curve over a
number field, and let $s:G_k\to\pi_1(V)$ be a continuous section.  For a
finite reduced divisor $D\subset V$, put $W=V-D$ and let
$\pi_1(W)^{\mathrm{cc},V}$ denote the maximal quotient of $\pi_1(W)$
whose kernel over $\pi_1(V)$ is abelian and has trivial conjugation
action by $\pi_1(V_{\bar k})$.  Then $s$ lifts continuously to
$\pi_1(W)^{\mathrm{cc},V}$.

At this fixed $V$, the lifts can be chosen compatibly for all $D$,
giving a continuous lift to
\[
 G_V^{\mathrm{cc}}:=
 \varprojlim_{D\subset V\text{ finite reduced}}
          \pi_1(V-D)^{\mathrm{cc},V}.
\]
The assertions hold after every finite extension of $k$ and on every
finite étale neighbourhood of the section.  In particular, on each
such neighbourhood the section annihilates the arithmetic Chern
classes of all line bundles.  No compatibility between the choices
on different finite étale neighbourhoods is asserted here.
\end{claim}
\begin{proof}
Write
\[
 T_D=\operatorname{Res}_{D/k}\mathbf G_m
     =\prod_{d\in|D|}\operatorname{Res}_{k(d)/k}\mathbf G_m,
 \qquad
 T(T_D)=\prod_{d\in|D|}
    \operatorname{Ind}_{G_{k(d)}}^{G_k}\widehat{\mathbf Z}(1).
\]
Let $\Delta\subset V\times_k V$ be the diagonal.  The torus torsor
\[
 F_D=\operatorname{Res}_{D\times_k V/V}
   \operatorname{Isom}
   \bigl(\mathcal O_{V\times V}(\Delta)|_{D\times V},
         \mathcal O_{D\times V}\bigr)\longrightarrow V
\]
is the construction of Borne--Emsalem--Stix
\cite[Definition 2.1]{D_BorneEmsalemStix2015}.  On $W$ the diagonal
does not meet $D\times W$, so its defining section gives a canonical
trivialization and a map $W\to F_D$ over $V$.

A hyperbolic curve in characteristic zero is an étale $K(\pi,1)$;
in particular $\pi_2^{\mathrm{et}}(V)=0$.  The torsor fibration and
its geometric action therefore give an exact sequence
\[
 1\longrightarrow T(T_D)\longrightarrow\pi_1(F_D)
   \longrightarrow\pi_1(V)\longrightarrow1,
 \tag{D.23.1}
\]
where the action on its kernel factors through $G_k$.
The canonical map from $W$ identifies
\[
 \pi_1(W)^{\mathrm{cc},V}\simeq\pi_1(F_D).
\]
These are precisely Propositions 3.1, 3.4 and 3.5 of
\cite{D_BorneEmsalemStix2015}.  Although the later applications in
that paper specialize to proper curves, these three statements and
the torsor construction apply to the present $V$.  Proposition 4.4
identifies the extension class with $c_1(F_D)$.  Thus the obstruction
to lifting $s$ is
\[
 c_D=s^*c_1(F_D)\in H^2\bigl(k,T(T_D)\bigr),
 \tag{D.23.2}
\]
with continuous cohomology of the inverse system of torsion modules;
its vanishing is equivalent to a continuous lift, as in
\cite[Lemma 4.5]{D_BorneEmsalemStix2015}.

We first prove that (D.23.2) vanishes at every place of $k$.  At a
finite place $v$, D.6 and D.19 say that $s_v$ is a point section or a
cuspidal section.  If it comes from $x_v\in V(k_v)$, the fibre
$F_{D,x_v}$ is a torsor under the quasi-trivial torus $T_D$.
Hilbert 90 gives $H^1(k_v,T_D)=0$, so that this fibre has a $k_v$-point.
That point provides the required local lift.

For a cusp $b\in(\overline V-V)(k_v)$, carry out the same diagonal
construction over the smooth completion $\overline V$.  It extends
$F_D$ and has a canonical section over $\overline V-D$, since the
diagonal is disjoint from $D\times(\overline V-D)$.  Restricting this
section to the punctured henselian neighbourhood of $b$ lifts its
whole decomposition group to $\pi_1(F_D)$.  Hence it lifts the given
cuspidal section, including its actual tangential cocycle.  No
rationality assumption on that cocycle is used.  At a real place,
Mochizuki's real section theorem identifies $s_v$ with a component
of $V(k_v)$, hence with the section of a point in that component
\cite[Corollary 3.13]{D_Mochizuki2003Real}; the point argument applies.
At a complex place the Galois group is trivial.  These arguments
show $c_D|_{k_v}=0$ for all $v$.

Here the local-to-global step has no inverse-limit ambiguity.  For a
number field $E$ and $n\mid m$, the transition
\[
 H^1(E,\mu_m)=E^\times/E^{\times m}
   \longrightarrow H^1(E,\mu_n)=E^\times/E^{\times n}
\]
is the natural surjection.  On the cofinal sequence $n!$ its
$\varprojlim{}^1$ is zero.  The Kummer sequence and Shapiro's lemma
therefore give
\[
 H^2\bigl(k,T(T_D)\bigr)
 \simeq \prod_{d\in|D|}\varprojlim_n
                   \operatorname{Br}(k(d))[n].
 \tag{D.23.3}
\]
Equivalently, this is Proposition 4.2 of
\cite{D_BorneEmsalemStix2015} for a quasi-trivial torus.
Under the local Shapiro isomorphisms, localization at all places of
$k$ includes every completion of every $k(d)$.  Brauer--Hasse--Noether
injectivity
\[
 \operatorname{Br}(E)\hookrightarrow
          \bigoplus_w\operatorname{Br}(E_w)
\]
for each of these number fields
\cite[Chapter VIII, Theorem 4.2]{D_MilneCFT2020}
annihilates every finite-level component of $c_D$.  Equation
(D.23.3) then gives $c_D=0$, proving existence of the lift.

For compatibility at the fixed base curve, choose once and for all
a geometric generic base point.  For disjoint reduced divisors one
has canonically
\[
 F_{D_1\amalg D_2}=F_{D_1}\times_V F_{D_2}.
\]
The exact sequences (D.23.1), applied also to the product torsor,
show that the natural map
\[
 \pi_1(F_{D_1\amalg D_2})
  \longrightarrow
 \pi_1(F_{D_1})\times_{\pi_1(V)}\pi_1(F_{D_2})
\]
is an isomorphism: it is an isomorphism on both the kernel
$T(T_{D_1})\times T(T_{D_2})$ and the quotient $\pi_1(V)$.
Choose a lift of $s$ for each one-point divisor $\{d\}$.  The fibre
products of these lifts give compatible lifts for every finite $D$.
Their inverse limit is continuous because each coordinate is
continuous.  This establishes the second assertion without taking
an inverse limit of arbitrary infinite sets of choices.

The argument applies unchanged over every finite extension and to
every finite étale neighbourhood.  Finally, a line bundle on a
regular curve is represented by an integral combination of closed
points.  For a point $d$, the norm homomorphism
$T_d\to\mathbf G_m$ sends the torsor Chern class just considered to
the Chern class of $\mathcal O_V(d)$, up to the harmless sign from
the convention for the isomorphism torsor.  Its pullback along $s$
is zero.  Additivity proves the claimed vanishing for all line
bundles, uniformly on those neighbourhoods.
\end{proof}

\begin{claim}{D.24}{Refuted}\label{clm:D24}
The following compatibility inference is false: a surjective trace
map between permutation Tate modules induces a surjection on
$H^1$ over a number field; consequently, after splitting all cusp
residue fields, arbitrary lower central lifts can be lifted upstairs
and descended by taking invariants.  The failure already occurs for
the trace from $\mathbf Q(i)$ to $\mathbf Q$ and the Kummer class
of $3$, detected modulo $2$.
\end{claim}
\begin{proof}
Let $L=\mathbf Q(i)$ and consider the surjection of $G_{\mathbf Q}$-modules
\[
 M_1=\operatorname{Ind}_{G_L}^{G_{\mathbf Q}}
             \widehat{\mathbf Z}(1)
   \xrightarrow{\operatorname{tr}}
 M_2=\widehat{\mathbf Z}(1).
 \tag{D.24.1}
\]
After restriction to $G_L$, this adds the two copies of
$\widehat{\mathbf Z}(1)$.  Shapiro, Kummer theory and the vanishing
of $\varprojlim{}^1$ for the finite groups $H^0(L,\mu_n)$ identify
its map on $H^1$ with
\[
 \widehat{L^\times}\xrightarrow{N_{L/\mathbf Q}}
           \widehat{\mathbf Q^\times},
 \qquad
 \widehat{E^\times}:=\varprojlim_n E^\times/E^{\times n}.
 \tag{D.24.2}
\]
This is the usual norm, rather than an inclusion or a choice of
one coordinate: the trace in (D.24.1) is the corestriction map under
Shapiro, and the Kummer connecting maps commute with corestriction.

These modules occur in an actual neighbourhood problem.  Take
\[
 V=\mathbf P^1_{\mathbf Q}-\{0,1,\infty\},\qquad
 V'=\mathbf P^1_{\mathbf Q}-\{0,1,-1,\infty\},\qquad
 t=z^2:V'\longrightarrow V.
\]
This is a finite étale cover.  The point section at $t=4$ lifts to
the point section at $z=2$.  Delete the lower point $d=-1$; its
inverse image in $V'$ is the degree-two divisor
$D'=\operatorname{Spec}\mathbf Q(i)$.  The lower and upper
cuspidally central kernels are respectively $M_2$ and $M_1$.
The cover is unramified at both geometric points over $d$, so each
upstairs inertia generator maps identically to the downstairs
generator; the induced kernel map is (D.24.1).  The actual point
sections at $4$ and $2$, which avoid these divisors, give compatible
origins for the two torsors of central lifts.  Twisting the lower
origin by $\kappa(3)$ cannot be lifted to an upper central lift of
the specified section at $z=2$, by the following norm obstruction.
Thus the trace failure is realized by the same finite étale
neighbourhood construction whose compatibility is at issue.

The class of $3$ is not in the image of (D.24.2).  Indeed, $3$ is
inert and unramified in $L$, and
\[
 v_3\bigl(N_{L/\mathbf Q}(a)\bigr)
       =2v_{(3)}(a)\qquad(a\in L^\times).
\]
Thus every norm has even $3$-adic valuation, even after multiplying
by a rational square, whereas $v_3(3)=1$.  This proves already that
\[
 L^\times/L^{\times2}\longrightarrow
 \mathbf Q^\times/\mathbf Q^{\times2}
\]
misses $[3]$.  Any hypothetical preimage in (D.24.2) would give a
preimage here, a contradiction.  An elementary certificate uses
$N(a+bi)=a^2+b^2$: the only solution of
$a^2+b^2=0$ in $\mathbf F_3^2$ is $(0,0)$, so the valuation of a
nonzero rational sum of two squares is even.  The complete
nine-pair check and the two-coset trace are recorded by
\path{code/D_norm_compatibility_certificate.py} and
\path{audit/D_norm_compatibility_certificate.json}.

Put $M=\ker(\operatorname{tr})$.  The long exact sequence also
shows that $H^2(\mathbf Q,M)\to H^2(\mathbf Q,M_1)$ need not be
injective: the connecting image of $[3]$ is nonzero and maps to
zero.  This is the precise failure in the older descent argument.

For source accuracy, the split-fibre surjectivity in Sa\"idi's
2010 preprint, printed p.37, is valid: a sum of copies of the
identity has an equivariant coordinate section when all fibres
are rational.  The failure is the subsequent assertion that
surjectivity survives taking
$\operatorname{Gal}(k_S/k)$-invariants, on pp.37--38, and the
alternative claim of surjectivity of
$H^1(G_k,M_1)\to H^1(G_k,M_2)$ on p.38
\cite[proof of Theorem 2.6(ii)]{D_Saidi2010}.
The map in that source sends each upstairs inertia summand
identically to its downstairs summand, and is exactly a trace of
permutation modules.  In the present example,
$H^0(G_L,M_1)=H^0(G_L,M_2)=0$, since a number field has no
nonzero cyclotomic Tate invariant.  Inflation--restriction thus
identifies the invariant $H^1$ groups over $L$ with the corresponding
$H^1$ groups over $\mathbf Q$.  The norm counterexample therefore
refutes the asserted invariant-surjectivity itself.

There is a separate, explicit finiteness issue in passing through
all neighbourhoods.  Sa\"idi's later condition $(H_\Sigma)$ requires
$H^1(G_k,A)$ to be finite for every finite $G_k$-module whose order
is supported on $\Sigma$; his inverse-limit argument in Lemma 2.3.4
uses nonempty finite sets of splittings
\cite[Definition 2.3.3 and Lemma 2.3.4]{D_Saidi2019}.
No number field satisfies this condition for a nonempty $\Sigma$:
for $\ell\in\Sigma$, the Kummer group
$H^1(k,\mu_\ell)=k^\times/k^{\times\ell}$ is infinite, as follows
by choosing successively elements with independent prime-valuation
coordinates using weak approximation.  Moreover,
$\widehat{\mathbf Q^\times}$ with its inverse-limit topology is
not compact, since its quotient
$\mathbf Q^\times/\mathbf Q^{\times2}$ is infinite and discrete.

D.23 consequently proves finite-rank and fixed-base central
lifting, while this argument prevents the specified automatic
descent or compactness step from supplying compatible lifts through
all finite étale neighbourhoods.  It does not prove that such a
compatible system fails to exist, and it does not refute the
cuspidalisation theorem under additional hypotheses, or D1--D4.
\end{proof}
